\begin{lemma}[Estimates for the small restricted source]
Let $B,A$ be a finite blow-up instance with candidacy system, host and
candidacy lower-$(\varepsilon,d)$-super-regularity, and host-size window
$(1-\varepsilon)n\leq |V_i|\leq(1+\varepsilon)n$. Suppose $\varepsilon\leq a$
and $(d/2)n\leq(d-\varepsilon)|V_i|$ for every $i$.
Then the host is lower-$(a,d/2)$-regular, its sizes lie in the window
$(1-a)n\leq |V_i|\leq(1+a)n$, and every original source vertex has at
least $(d/2)n$ original candidates.
\end{lemma}
\begin{proof}
For an edge $ij$ of the reduced graph, \noderef{InstanceLowerSuperRegular}
gives the original lower-super-regular host pair. Its first clause in
\noderef{LowerSuperRegularPair} is lower regularity.
Take subsets of its two parts of sizes strictly greater than $a|V_i|$
and $a|V_j|$. Since $\varepsilon\leq a$, these sizes are also strictly
greater than the original $\varepsilon$ thresholds. Thus
\noderef{LowerRegularPair} gives density at least $d-\varepsilon$.
Because $d/2\leq d$ and $\varepsilon\leq a$, monotonicity of truncated
subtraction gives $d/2-a\leq d-\varepsilon$.
This is precisely \noderef{InstanceLowerRegular} at $(a,d/2)$.
The size-window conclusions follow from $1-a\leq1-\varepsilon$
and $1+\varepsilon\leq1+a$, multiplying by nonnegative $n$.

For any vertex $x$, put $i=c_H(x)$. In the sum-vertex graph defining
\noderef{CandidacyLowerSuperRegular}, the vertex $\mathrm{inl}(x)$ lies
in the left $i$th part. The left-degree clause of
\noderef{LowerSuperRegularPair} gives at least
$(d-\varepsilon)|V_i|$ neighbors in the right part. That right part is
bijective with $V_i$ by $v\mapsto\mathrm{inr}(v)$.
For opposite tags, expanding the defining relation of the graph shows
that adjacency of $\mathrm{inl}(x)$ and $\mathrm{inr}(v)$ is exactly
$A(x,v)$: the forward orientation is this relation, the reverse
orientation cannot have the stipulated tags, and the endpoints are distinct.
Moreover \noderef{CandidacySystem} forces every allowed $v$ into $V_i$.
Thus the filtered neighbor set is bijective with the entire original
candidate set of $x$. Its cardinality is at least
$(d-\varepsilon)|V_i|\geq(d/2)n$, as required.
\end{proof}
